% =====================================================================
%  Brand style shared by every quote and proposal: packages, colours,
%  layout and reusable section blocks.
%
%  This file holds NO company data. Names, tax id, phones, address and
%  image file names come from a separate company file that must be
%  loaded first (see company.example.tex).
%
%  Usage, from a document:
%
%      \documentclass[11pt,a4paper]{article}
%      \newif\ifdraft
%      \drafttrue                      % optional; the default is final
%      \newcommand{\CoLang}{en}        % optional; the default is Spanish
%      \newcommand{\CoAssetsDir}{/path/to/branding/assets/}
%      \input{/path/to/company}
%      \input{/path/to/colcastar_style}
%      \begin{document}
%      \letterhead
%      ...
%
%  \ifdraft, \CoLang and \CoAssetsDir must be set BEFORE this file is
%  loaded, because \TODO, the language and the image paths depend on them.
% =====================================================================

\ifdefined\CoLegalName\else
  \errmessage{Load the company data file before colcastar_style.tex}
\fi

% =====================================================================
%  1. Packages
% =====================================================================
% Language. The text is compared, not the meaning: \newcommand makes
% macros \long, and an \ifx against a plain \def never matches.
\usepackage{etoolbox}
\providecommand{\CoLang}{es}
\newif\ifcoEnglish
\ifdefstring{\CoLang}{en}{\coEnglishtrue}{\coEnglishfalse}

% A document that does not declare the draft switch is a final version.
\ifcsname ifdraft\endcsname\else
  \expandafter\newif\csname ifdraft\endcsname
\fi

% Fonts. pdfLaTeX uses the classic Helvetica clone; XeTeX and LuaTeX
% (tectonic) need a Unicode font, or characters such as the inverted
% question mark land on the wrong glyph.
\usepackage{iftex}
\ifPDFTeX
  \usepackage[T1]{fontenc}
  \usepackage[utf8]{inputenc}
  \usepackage[scaled=0.95]{helvet}
\else
  \usepackage{fontspec}
  \setsansfont{texgyreheros}[
    Extension=.otf,
    UprightFont=*-regular,
    BoldFont=*-bold,
    ItalicFont=*-italic,
    BoldItalicFont=*-bolditalic,
    Scale=0.95]
\fi
\ifcoEnglish
  \usepackage[english]{babel}
\else
  \usepackage[spanish,es-noshorthands]{babel}
\fi
\renewcommand{\familydefault}{\sfdefault}

\usepackage[a4paper,top=2cm,bottom=2.2cm,left=2cm,right=2cm]{geometry}
\usepackage[table]{xcolor}
\usepackage{graphicx}
\usepackage{tabularx}
\usepackage{array}
\usepackage{booktabs}
\usepackage{enumitem}
\usepackage{fancyhdr}
\usepackage{titlesec}
\usepackage{parskip}
\usepackage[hidelinks]{hyperref}
\usepackage{needspace}

% Company images live in one folder. \CoAssetsDir is the path to it as
% seen from the document. A local assets/ folder is searched first, for
% images that belong to one document and are not brand material.
\providecommand{\CoAssetsDir}{../branding/assets/}
\edef\CoGraphicsPath{\noexpand\graphicspath{{assets/}{\CoAssetsDir}}}
\CoGraphicsPath

% =====================================================================
%  2. Brand colours
% =====================================================================
\definecolor{brandnavy}{HTML}{0F172A}
\definecolor{brandgold}{HTML}{F59E0B}
\definecolor{brandslate}{HTML}{475569}
\definecolor{brandmist}{HTML}{F1F5F9}
\definecolor{brandline}{HTML}{CBD5E1}
\definecolor{todomark}{HTML}{C2410C}

\hypersetup{
  colorlinks=true, urlcolor=brandnavy, linkcolor=brandnavy,
  pdfauthor={\CoLegalName}
}

% =====================================================================
%  3. Draft marks
% =====================================================================
\ifdraft
  \newcommand{\TODO}[1]{{\leavevmode\bfseries\color{todomark}[#1]}}
\else
  \newcommand{\TODO}[1]{[#1]}
\fi

% =====================================================================
%  4. Layout
%  ---------------------------------------------------------------------
%  A p-column cell whose content starts with \color begins its vertical
%  list with a whatsit, the \vtop gets zero height and the content drops
%  below the row baseline. \leavevmode forces a first box and realigns.
%  It goes before every \color that opens a cell.
% =====================================================================
\setlength{\arrayrulewidth}{0.4pt}
\renewcommand{\arraystretch}{1.3}
\arrayrulecolor{brandline}

\titleformat{\section}
  {\color{brandnavy}\normalfont\large\bfseries}
  {}{0pt}{}[\vspace{-0.55em}\color{brandgold}\rule{\linewidth}{2pt}]
\titlespacing*{\section}{0pt}{1.5em}{0.7em}

\pagestyle{fancy}
\fancyhf{}
\renewcommand{\headrulewidth}{0pt}
\renewcommand{\footrulewidth}{0.4pt}
\fancyfoot[L]{\footnotesize\color{brandslate}\CoLegalName\quad RUC \CoRuc}
\fancyfoot[R]{\footnotesize\color{brandslate}\ifcoEnglish Page\else Página\fi\ \thepage}

% An image that does not break the build when the file is missing.
% Pass ONLY the file name. It looks in the document's own assets/ first
% and then in \CoAssetsDir. \IfFileExists does NOT honour \graphicspath,
% so the paths are tried by hand.
\newcommand{\safeimg}[2]{%
  \IfFileExists{assets/#1}%
    {\includegraphics[height=#2,keepaspectratio]{assets/#1}}%
    {\IfFileExists{\CoAssetsDir#1}%
      {\includegraphics[height=#2,keepaspectratio]{\CoAssetsDir#1}}%
      {\fbox{\parbox[c][#2][c]{4cm}{\centering\footnotesize\color{brandslate}missing image:\\\coverbatim{#1}}}}}}

% Prints a file name or path literally. Names with an underscore or another
% TeX-special character would otherwise break the build in text mode.
\newcommand{\coverbatim}[1]{\texttt{\expandafter\detokenize\expandafter{#1}}}

% Runs #2 when image #1 is found where \safeimg and \cologo look for it
% (the document's own assets/ first, then \CoAssetsDir), and #3 otherwise.
\newcommand{\coimgstatus}[3]{%
  \IfFileExists{assets/#1}{#2}{\IfFileExists{\CoAssetsDir#1}{#2}{#3}}}

% The logo, with the same fallback, sized by width for the letterhead.
\newcommand{\cologo}{%
  \IfFileExists{assets/\CoLogoFile}%
    {\includegraphics[width=3.1cm]{assets/\CoLogoFile}}%
    {\IfFileExists{\CoAssetsDir\CoLogoFile}%
      {\includegraphics[width=3.1cm]{\CoAssetsDir\CoLogoFile}}%
      {\fbox{\parbox[c][1.6cm][c]{2.7cm}{\centering\footnotesize\color{brandslate}logo}}}}}

% Letterhead. Called once, right after the document opens.
\newcommand{\letterhead}{%
  \noindent
  \begin{tabularx}{\linewidth}{@{}p{3.2cm}X@{}}
    \raisebox{-0.35\height}{\cologo} &
    \raggedleft
    {\Large\bfseries\color{brandnavy}\CoName}\\[0.15em]
    {\footnotesize\color{brandslate}\CoLegalName\ --- RUC \CoRuc}\\
    {\footnotesize\color{brandslate}\CoCity}\\
    {\footnotesize\color{brandslate}\href{mailto:\CoEmail}{\CoEmail} \textbullet\ \CoPhone}\\
    {\footnotesize\color{brandslate}\href{\CoWebUrl}{\CoWeb}}
  \end{tabularx}
  \par\vspace{0.6em}
  {\color{brandgold}\rule{\linewidth}{3pt}}\par}

% Closing signature.
\newcommand{\signature}[1]{%
  \par\vspace{1.1em}\noindent
  \ifcoEnglish Kind regards,\else Un cordial saludo,\fi\\[0.6em]
  {\bfseries\color{brandnavy}#1}\\
  {\footnotesize\color{brandslate}\CoLegalName}\\
  {\footnotesize\color{brandslate}\href{mailto:\CoEmail}{\CoEmail} \textbullet\ \CoPhone\ \textbullet\ \href{\CoWebUrl}{\CoWeb}}\par}

% Section title. The optional argument is how many lines must be free at
% the bottom of the page before starting, so the title is not orphaned.
% Raise it when what follows is a block that cannot be split.
\newcommand{\sect}[2][6]{\Needspace*{#1\baselineskip}\section*{#2}}

% Day header inside an itinerary.
\newcommand{\dayhead}[2]{%
  \par\vspace{1.0em}\noindent
  \begin{tabularx}{\linewidth}{@{}>{\columncolor{brandnavy}}p{3.3cm}>{\columncolor{brandmist}}X@{}}
    \leavevmode\color{white}\bfseries #1 & \leavevmode\color{brandnavy}\bfseries #2 \\
  \end{tabularx}\par\vspace{0.45em}}

% Option header, when a proposal offers more than one programme.
\newcommand{\opthead}[2]{%
  \par\vspace{1.1em}\noindent
  \begin{tabularx}{\linewidth}{@{}>{\columncolor{brandgold}}p{3.3cm}>{\columncolor{brandnavy}}X@{}}
    \leavevmode\color{brandnavy}\bfseries #1 & \leavevmode\color{white}\bfseries #2 \\
  \end{tabularx}\par\vspace{0.5em}}

% Day footer: meals and accommodation.
\newcommand{\daymeta}[2]{%
  \par\vspace{0.2em}%
  {\footnotesize\color{brandslate}%
  \textbf{\ifcoEnglish Meals:\else Comidas:\fi} #1 \quad\textbullet\quad
  \textbf{\ifcoEnglish Accommodation:\else Alojamiento:\fi} #2}\par}

% Highlighted note, with the gold bar at the side.
\newcommand{\notebox}[1]{%
  \par\vspace{0.5em}\noindent
  \textcolor{brandgold}{\rule{2pt}{\dimexpr\baselineskip\relax}}\hspace{0.5em}%
  \begin{minipage}{\dimexpr\linewidth-1.2em\relax}
    {\footnotesize #1}
  \end{minipage}\par\vspace{0.3em}}

% =====================================================================
%  5. Reusable section blocks
%  ---------------------------------------------------------------------
%  These sections are the same in every proposal because they describe
%  the company, not the client. They update themselves when the company
%  file changes.
% =====================================================================

% Diagnostics: prints where the paths resolved and whether each image is
% found. Put it on a draft page when an image does not show up.
\newcommand{\brandingdiag}{%
  \par{\footnotesize\color{brandslate}
  \textbf{Resolved paths.}   Image folder: \coverbatim{\CoAssetsDir}\\
  \coimgstatus{\CoLogoFile}{Logo: found}{\textbf{Logo: NOT found}}\\
  \coimgstatus{\CoConstanciaFile}{Certificate: found}{\textbf{Certificate: NOT found}}\\
  \coimgstatus{\CoVehicleFileA}{Vehicle photo A: found}{\textbf{Vehicle photo A: NOT found}}\\
  \coimgstatus{\CoVehicleFileB}{Vehicle photo B: found}{\textbf{Vehicle photo B: NOT found}}
  \par}}

% --- Your vehicle, driver and guide -----------------------------------
% #1 = the "for your group" line, which depends on the client.
% #2 = the guide line, which changes with the language and the days.
\newcommand{\vehicletable}[2]{%
  \noindent\makebox[\linewidth][c]{%
    \safeimg{\CoVehicleFileA}{5.0cm}\hspace{0.7em}%
    \safeimg{\CoVehicleFileB}{5.0cm}}
  \par\vspace{0.7em}\noindent
  \begin{tabularx}{\linewidth}{@{}>{\leavevmode\bfseries\color{brandnavy}}l@{\hspace{1.4em}}>{\raggedright\arraybackslash}X@{}}
    \ifcoEnglish Vehicle\else Vehículo\fi & \ifcoEnglish\CoVehicleEn\else\CoVehicle\fi \\
    \ifcoEnglish For your group\else Para su grupo\fi & #1 \\
    \ifcoEnglish Permits\else Permisos\fi & \ifcoEnglish\CoVehiclePermitEn\else\CoVehiclePermit\fi \\
    \ifcoEnglish Driver\else Conductor\fi &
      \ifcoEnglish Extensive experience on the Colca routes, with local knowledge of the valley%
      \else Amplia experiencia en las rutas del Colca, con conocimiento local del valle\fi \\
    \ifcoEnglish Guide\else Guía\fi & #2 \\
  \end{tabularx}\par}

% --- Our credentials ---------------------------------------------------
% No arguments: everything comes from the company file.
\newcommand{\credentials}{%
  \noindent
  \begin{minipage}[t]{6.4cm}
    \vspace{0pt}%
    \safeimg{\CoConstanciaFile}{7.6cm}
  \end{minipage}\hfill
  \begin{minipage}[t]{\dimexpr\linewidth-6.9cm\relax}
    \vspace{0pt}%
    \raggedright
    \ifcoEnglish
      We are a travel agency and tour operator formally registered with the
      \textbf{Regional Government of Arequipa}, through its Regional Board of Foreign
      Trade and Tourism (GRCET).

      \vspace{0.5em}
      The certificate shown here is our operating licence. It certifies that
      \textbf{\CoLegalName}, RUC \CoRuc, has met the requirements of Supreme Decree
      N.\textordmasculine{} 026-2004-MINCETUR and is classified as a
      \textbf{Tour Operator}.

      \vspace{0.5em}
      {\footnotesize\color{brandslate}\CoConstancia, \CoConstanciaFechaEn.}
    \else
      Somos una agencia de viajes y operador de turismo formalmente registrada ante el
      \textbf{Gobierno Regional de Arequipa}, a través de su Gerencia Regional de Comercio
      Exterior y Turismo (GRCET).

      \vspace{0.5em}
      La constancia que acompaña a este texto es nuestra autorización de funcionamiento.
      Acredita que \textbf{\CoLegalName}, RUC \CoRuc, ha cumplido con los requisitos del
      Decreto Supremo N.\textordmasculine{} 026-2004-MINCETUR y está clasificada como
      \textbf{Operador de Turismo}.

      \vspace{0.5em}
      {\footnotesize\color{brandslate}\CoConstancia, \CoConstanciaFecha.}
    \fi
  \end{minipage}
  \par\vspace{0.8em}\noindent
  \begin{tabularx}{\linewidth}{@{}>{\leavevmode\bfseries\color{brandnavy}}l@{\hspace{1.4em}}>{\raggedright\arraybackslash}X@{}}
    \ifcoEnglish Registered name\else Razón social\fi & \CoLegalName \\
    \ifcoEnglish Trading name\else Nombre comercial\fi & \CoTradeName \\
    RUC & \CoRuc \\
    \ifcoEnglish Activity\else Actividad\fi & \ifcoEnglish\CoCiiuEn\else\CoCiiu\fi \\
    \ifcoEnglish Registered address\else Domicilio fiscal\fi & \CoAddress \\
    \ifcoEnglish Founder\else Fundador\fi & \ifcoEnglish\CoFounderEn\else\CoFounder\fi \\
  \end{tabularx}\par}

% --- To confirm ---------------------------------------------------------
% #1 = what we need the client to answer. It depends on the message they
%      sent, so the quote passes it in.
\newcommand{\confirmblock}[1]{%
  \noindent #1
  \par\vspace{0.6em}\noindent
  \begin{tabularx}{\linewidth}{@{}>{\leavevmode\bfseries\color{brandnavy}}l@{\hspace{1.4em}}>{\raggedright\arraybackslash}X@{}}
    WhatsApp & \CoWhatsApp\ \ifcoEnglish(official)\else(oficial)\fi \\
    \ifcoEnglish Emergencies\else Emergencias\fi & \CoEmergency \\
    \ifcoEnglish Email\else Correo\fi & \href{mailto:\CoEmail}{\CoEmail} \\
    \ifcoEnglish Website\else Web\fi & \href{\CoWebUrl}{\CoWeb} \\
    \ifcoEnglish Payment\else Pago\fi & \ifcoEnglish\CoPaymentEn\else\CoPayment\fi \\
  \end{tabularx}\par}
